\subsection{Neural solution and action-domain sensitivity}\label{sec:ndu-neural}
We now apply multilayer NBO to the same two-state economy, with the original cardinal utility, correlation, reflection rule, and stopping payoff unchanged. At each backward time level a two-hidden-layer, width-thirty-two tanh critic approximates its own Bellman continuation; a multilayer categorical actor selects among feasible actions. The reference solution is computed separately and is not used as a training target. The actor is fitted before policy evaluation at that time level, with the next-period critic fixed. The direct comparator eliminates the actor by enumeration of the same action set. Both methods have a cap of 160 L-BFGS inner iterations per time level, divided 100/60 between critic and actor for NBO.

An initial unweighted classification implementation produced a large policy loss. Absorbing wealth-face labels, for which every action has identical value, contaminated the policy fit; classification error also ignored the economic cost of a mistaken low-consumption action. The retained implementation excludes those faces from actor training and minimizes expected normalized Bellman action loss, with a small cross-entropy term. The pilot source, weights, and failed output are preserved. This change repairs an implementation failure rather than redefining the economic objective: the deployed action is still evaluated against \emph{all} finite feasible alternatives.

For the principal comparison, $(n_u,n_x,N_t)=(17,25,20)$ and $k=1$. The expanded action set contains 175 triples, with $c/X\in\{.02,.05,.10,.20,.30\}$, $p\in\{0,.375,.75,1.125,1.5\}$, and $\theta\in\{-.45,-.30,-.15,0,.15,.30,.45\}$. Every time--state--action combination enters the error accounts in \eqref{eq:accounts-r6}. Independently solving the fixed learned policy gives its actual regret relative to the finite reference. Midpoint value errors and derivative comparisons are additional diagnostics, not continuous-domain certificates.
\begin{table}[t]\centering
\caption{Neural preference adjustment on the same finite economy}\label{tab:ndu-r6}
\small\begin{tabular}{lrrrrr}
\toprule
Method & Seed & Value error & Policy regret & Bound $R_0$ & Midpoint RMS\\\midrule
NBO & 11 & 0.1586 & 0.1308 & 2.7443 & 0.0433\\
NBO & 29 & 0.1546 & 0.0790 & 2.7493 & 0.0407\\
NBO & 47 & 0.1523 & 0.0364 & 2.9394 & 0.0454\\
Direct Bellman & 11 & 0.1114 & 0.0261 & 1.7752 & 0.0476\\
Direct Bellman & 29 & 0.1104 & 0.0260 & 1.7443 & 0.0402\\
Direct Bellman & 47 & 0.1301 & 0.0268 & 1.8383 & 0.0466\\
\bottomrule\end{tabular}
\par\smallskip\footnotesize The first three error columns use every state and time of the $17\times25\times21$ grid with 175 actions; $R_0$ is the time-zero certificate, while the first two columns maximize their independently computed errors over all times. Midpoints compare the time-zero neural critic with the finite-reference interpolant. Continuous-action and controlled-diffusion discretization errors are not included. The accuracy targets were value error $.2$ and policy regret $.1$.
\end{table}

The three-seed record contains a failed NBO accuracy target and reports it without exclusion. The direct comparator is more accurate at this budget. A categorical actor therefore has a measurable approximation cost even when the critic class is identical. Neural representation alone is not a reason to replace an inexpensive two-state grid solver; the grid here supplies an unusually complete audit of the neural computation.

Action bounds are part of the economy, not innocuous numerical constants. The sensitivity study reproduces the referee's action-set comparison and extends it to 1,053 actions in the same expanded box, a wider consumption--portfolio--adjustment box, a wider wealth domain, and two finer space--time grids, for both $k=1$ and $k=5$. The table in the supplement records lower- and upper-bound frequencies for every control. Consumption and portfolio bounds remain frequently active after the first expansion. Thus state-grid convergence of the old 27-action model does not establish convergence to an unconstrained control optimum. The wider sets describe different constrained economies; refinement within a fixed box is a different comparison. No continuous-action coverage or controlled-diffusion discretization error is set to zero in the finite certificate. This preserves the economic application while making its constraint dependence observable.
